\begin{table}[H]\centering\small
\caption{Caractéristiques des surfaces et du contact sous la pression nominale du TFE (\SI{1.65}{MPa}).
$m_2$ : pente quadratique moyenne ; aire : fraction de la surface en contact ; $p_{99}$ : 99\textsuperscript{e}
centile de la pression locale dans le contact (pour les rainures : maximum, qui dépend de la
résolution à cause de la singularité d'arête).}\label{tab:surfaces}
\begin{tabular}{lccccccc}
\toprule
Surface & MPD & ETD & $S_q$ & $S_{sk}$ & $m_2$ & aire & $p_{99}$\\
 & (mm) & (mm) & (mm) & & & (\%) & (MPa)\\
\midrule
rainures FAA, arêtes vives & --- & --- & --- & --- & --- & 84 & 7,2 (max)\\
rainures FAA, arêtes $r$ = 1 mm & --- & --- & --- & --- & --- & 82 & 4,3 (max)\\
aléatoire, MPD 0,6 mm & 0,60 & 0,68 & 0,32 & $-$0,03 & 0,050 & 76 & 5,5\\
aléatoire, MPD 1,0 mm & 1,00 & 1,00 & 0,53 & $-$0,03 & 0,140 & 55 & 8,2\\
aléatoire, MPD 1,5 mm & 1,50 & 1,40 & 0,80 & $-$0,03 & 0,315 & 40 & 11,7\\
négative, MPD 1,0 mm & 1,00 & 1,00 & 0,74 & $-$1,28 & 0,289 & 52 & 7,2\\
positive, MPD 1,0 mm & 1,00 & 1,00 & 0,41 & 1,28 & 0,089 & 70 & 9,5\\
MPD 1,0 ; $E_r$ = 5 MPa & 1,00 & 1,00 & 0,53 & $-$0,03 & 0,140 & 83 & 4,8\\
MPD 1,0 ; $E_r$ = 20 MPa & 1,00 & 1,00 & 0,53 & $-$0,03 & 0,140 & 31 & 15,2\\
\bottomrule
\end{tabular}
\end{table}
